% main.tex
% gs -dNOPAUSE -dBATCH -sDEVICE=png16m -sOutputFile=out%03d.png -r1000x1000 main.pdf
% ffmpeg -framerate 30 -i out%03d.png out.mp4
% https://tex.stackexchange.com/a/748667
\unprotect
    %%% Defined similarly to \obeyMPboxdepth, \ignoreMPboxdepth,
    %%% \obeyMPboxorigin, and \normalMPboxdepth
    \permanent\protected\def\strutMPboxdepth{%
        \let\meta_relocate_box\meta_strut_box_depth%
    }

    \protected\def\meta_strut_box_depth{%
        \setbox\b_meta_graphic\vpack{%
            \box\b_meta_graphic%
            %%% Add some depth to the box without affecting its size
            \nohrule height-\openstrutdepth width 0pt depth \openstrutdepth\relax%
        }
    }

    %%% The pagecolumns output routine alredy adds a \vfil elsewhere, so the one
    %%% here isn't needed and disrupts the depth added by \strutMPboxdepth.
    \let\page_otr_fill_page_unchecked\relax
\protect

%%% Comment the line below to confirm that an overfull vbox is generated;
%%% uncomment it to confirm that the overfull vbox message goes away.
\strutMPboxdepth
\setupbodyfont[pagella]


\definepapersize
  [youtubeHD]
  [width=1920bp,height=1080bp]
\setuppapersize[youtubeHD][youtubeHD]
\setuplayout[
    location=middle,
    width=middle,
    height=middle,
    topspace=0pt,
    backspace=0pt,
    header=0pt,
    footer=0pt
]
\defineframed[myframed][
    width=800pt,
    align=normal,
    offset=1cm,
    frame=off,
    indenting={yes,medium}
]
\setupbodyfont[modern,24pt]
\ctxloadluafile{parametric.lua}
\processMPfigurefile{register_mp_cmd.mp}

\startMPinclusions
def render_theorem_box =
numeric w, h;
w := CurrentWidth;
h := CurrentHeight;

path p;
p := (-w/2, -h/2) -- (w/2, -h/2) -- (w/2, h/2) -- (-w/2, h/2) -- cycle;
setbounds currentpicture to p;

picture theorem;
theorem := textext("\getbuffer[BufA]");

pair center_anchor, topleft_anchor;
center_anchor := center theorem;
topleft_anchor := ulcorner theorem;

enddef;


def jasper_thin =
    withpen pencircle scaled .1cm
enddef ;

def jasper_thick =
    withpen pencircle scaled .2cm
enddef ;

def jasper_lightgray =
    withcolor (0.6)
enddef ;

def jasper_lightgraypenthin =
    jasper_lightgray withpen pencircle scaled .1cm
enddef ;

def jasper_lightgraypenthick =
    jasper_lightgray withpen pencircle scaled .2cm
enddef ;


def jasper_darkgray =
    withcolor (0.4)
enddef ;

def jasper_darkgraypenthin =
    jasper_darkgray withpen pencircle scaled .15cm
enddef ;

def jasper_darkgraypenthick =
    jasper_darkgray withpen pencircle scaled .3cm
enddef ;
\stopMPinclusions




\starttext
    \startbuffer[BufA]
        \startframed[myframed]
            {\bf AXIOM 1} Given any two points in space, one and 
            only one line can be drawn through these points.\par
            \color[white]{It follows from AXIOM 1 that if any two points of a 
            line coincide with two points of some other line, 
            then these lines coincide in all their points.
            (Otherwise it would be possible to draw two 
            distinct lines through a given pair of points).\par 
            Therefore, two distinct lines may intersect only at 
            one point.\par 
            Let us emphasize that AXIOM 1 asserts not only the 
            existence of a line through a pair of points but also 
            the fact that such a line can be drawn.\par 
            Thus AXIOM 1 implies the existence of an instrument 
            that allows us to draw a straight line through any 
            pair of points. Such an instrument is called a 
            straightedge.}
        \stopframed
    \stopbuffer
    % Display the theorem initially, in a prominent location.
    \dorecurse{7*30}{%
        \startMPpage
            render_theorem_box ;
            numeric t;
            t := 0;
            pair anchor;
            anchor := (1 - t) * center_anchor + t * topleft_anchor + t * (-2cm, -9cm);
            label(theorem, (0,3cm) - anchor);
        \stopMPpage%
    }
    % Translate the theorem to the top right corner.
    \dorecurse{2*30+1}{%
        \startMPpage
            render_theorem_box ;
            numeric t;
            t := (#1-1)/60;
            pair anchor;
            anchor := (1 - t) * center_anchor + t * topleft_anchor + t * (-2cm, -9cm);
            label(theorem, (0,3cm) - anchor);
        \stopMPpage%
    }
    % Pause once the theorem is in place.
    \dorecurse{3*30}{%
        \startMPpage
            render_theorem_box ;
            numeric t;
            t := 1;
            pair anchor;
            anchor := (1 - t) * center_anchor + t * topleft_anchor + t * (-2cm, -9cm);
            label(theorem, (0,3cm) - anchor);
        \stopMPpage%
    }
    \startMPinclusions
        def mytext =
            render_theorem_box ;
            numeric t;
            t := 1; 
            pair anchor;
            anchor := (1 - t) * center_anchor + t * topleft_anchor + t * (-2cm, -9cm);
            label(theorem, (0,3cm) - anchor);
        enddef ;
        def points_LO =
            jasper_point[
                fx="cx"
                ,fy="cy+5"
            ] ;
            jasper_point[
                fx="cx"
                ,fy="cy-5"
            ] ;
        enddef ;
    \stopMPinclusions%
    % draw two points
    \dorecurse{4*30}{%
        \startMPpage
            mytext ;
            points_LO ;
            jasper_render_segments ;
        \stopMPpage%
    }
    % draw a line through the points
    \dorecurse{3*30}{%
        \startMPpage
            mytext ;
            points_LO ;
            jasper_curve[
                ustart="0"
                ,ustop="1"
                ,usamples="2"
                ,fx="cx"
                ,fy="cy+7-14*u*#1/90"
            ] ;
            jasper_render_segments ;
        \stopMPpage%
    }
    % pause after drawing the line
    \startMPinclusions
        def lineone =
            points_LO ;
            jasper_curve[
                ustart="0"
                ,ustop="1"
                ,usamples="2"
                ,fx="cx"
                ,fy="cy+7-14*u"
            ] ;
        enddef ;
    \stopMPinclusions%
    \dorecurse{3*30}{%
        \startMPpage
            mytext ;
            lineone ;
            jasper_render_segments ;
        \stopMPpage%
    }
    \startbuffer[BufA]
        \startframed[myframed]
            {\bf AXIOM 1} Given any two points in space, one and 
            only one line can be drawn through these points.\par
            It follows from AXIOM 1 that if any two points of a 
            line coincide with two points of some other line, 
            then these lines coincide in all their points.
            \color[white]{(Otherwise it would be possible to draw two 
            distinct lines through a given pair of points).\par 
            Therefore, two distinct lines may intersect only at 
            one point.\par 
            Let us emphasize that AXIOM 1 asserts not only the 
            existence of a line through a pair of points but also 
            the fact that such a line can be drawn.\par 
            Thus AXIOM 1 implies the existence of an instrument 
            that allows us to draw a straight line through any 
            pair of points. Such an instrument is called a 
            straightedge.}
        \stopframed
    \stopbuffer
    % It follows from AXIOM 1 that if any two points of a 
    % line coincide with two points of some other line, 
    % then these lines coincide in all their points.
    \dorecurse{10*30}{%
        \startMPpage
            mytext ;
            lineone ;
            jasper_render_segments ;
        \stopMPpage%
    }
    % draw two points of second line
    \dorecurse{4*30}{%
        \startMPpage
            mytext ;
            lineone ;
            jasper_point[
                fx="cx"
                ,fy="cy+5"
                ,transformation="
                    matrix_multiply(
                        inverse(translate(cx,cy-13,0))
                        ,matrix_multiply(
                            zrotation(pi/6)
                            ,translate(cx,cy-13,0)
                        )
                    )
                "
            ] ;
            jasper_point[
                fx="cx"
                ,fy="cy-5"
                ,transformation="
                    matrix_multiply(
                        inverse(translate(cx,cy-13,0))
                        ,matrix_multiply(
                            zrotation(pi/6)
                            ,translate(cx,cy-13,0)
                        )
                    )
                "
            ] ;
            jasper_render_segments ;
        \stopMPpage%
    }
    % draw line through those points
    \dorecurse{3*30}{%
        \startMPpage
            mytext ;
            lineone ;
            jasper_point[
                fx="cx"
                ,fy="cy+5"
                ,transformation="
                    matrix_multiply(
                        inverse(translate(cx,cy-13,0))
                        ,matrix_multiply(
                            zrotation(pi/6)
                            ,translate(cx,cy-13,0)
                        )
                    )
                "
            ] ;
            jasper_point[
                fx="cx"
                ,fy="cy-5"
                ,transformation="
                    matrix_multiply(
                        inverse(translate(cx,cy-13,0))
                        ,matrix_multiply(
                            zrotation(pi/6)
                            ,translate(cx,cy-13,0)
                        )
                    )
                "
            ] ;
            jasper_curve[
                ustart="0"
                ,ustop="1"
                ,usamples="2"
                ,fx="cx"
                ,fy="cy+7-14*u*#1/90"
                ,transformation="
                    matrix_multiply(
                        inverse(translate(cx,cy-13,0))
                        ,matrix_multiply(
                            zrotation(pi/6)
                            ,translate(cx,cy-13,0)
                        )
                    )
                "
            ] ;
            jasper_render_segments ;
        \stopMPpage%
    }
    % pause after drawing line through new pair of points
    \dorecurse{4*30}{%
        \startMPpage
            mytext ;
            lineone ;
            jasper_point[
                fx="cx"
                ,fy="cy+5"
                ,transformation="
                    matrix_multiply(
                        inverse(translate(cx,cy-13,0))
                        ,matrix_multiply(
                            zrotation(pi/6)
                            ,translate(cx,cy-13,0)
                        )
                    )
                "
            ] ;
            jasper_point[
                fx="cx"
                ,fy="cy-5"
                ,transformation="
                    matrix_multiply(
                        inverse(translate(cx,cy-13,0))
                        ,matrix_multiply(
                            zrotation(pi/6)
                            ,translate(cx,cy-13,0)
                        )
                    )
                "
            ] ;
            jasper_curve[
                ustart="0"
                ,ustop="1"
                ,usamples="2"
                ,fx="cx"
                ,fy="cy+7-14*u"
                ,transformation="
                    matrix_multiply(
                        inverse(translate(cx,cy-13,0))
                        ,matrix_multiply(
                            zrotation(pi/6)
                            ,translate(cx,cy-13,0)
                        )
                    )
                "
            ] ;
            jasper_render_segments ;
        \stopMPpage%
    }
    % undo the transformation to coincide the lines
    \def\myT{
        matrix_multiply(
            inverse(translate(cx,cy-13,0))
            ,matrix_multiply(
                zrotation(pi/6)
                ,translate(cx,cy-13,0)
            )
        )
    }%
    \dorecurse{3*30}{%
        \startMPpage
            mytext ;
            lineone ;
            
            jasper_point[
                fx="cx"
                ,fy="cy+5"
                ,transformation="
                    matrix_multiply(
                        inverse(translate(cx,cy-13,0))
                        ,matrix_multiply(
                            zrotation(pi/6*(1-#1/90))
                            ,translate(cx,cy-13,0)
                        )
                    )
                "
            ] ;
            jasper_point[
                fx="cx"
                ,fy="cy-5"
                ,transformation="
                    matrix_multiply(
                        inverse(translate(cx,cy-13,0))
                        ,matrix_multiply(
                            zrotation(pi/6*(1-#1/90))
                            ,translate(cx,cy-13,0)
                        )
                    )
                "
            ] ;
            jasper_curve[
                ustart="0"
                ,ustop="1"
                ,usamples="2"
                ,fx="cx"
                ,fy="cy+7-14*u"
                ,transformation="
                    matrix_multiply(
                        inverse(translate(cx,cy-13,0))
                        ,matrix_multiply(
                            zrotation(pi/6*(1-#1/90))
                            ,translate(cx,cy-13,0)
                        )
                    )
                "
            ] ;
            jasper_render_segments ;
        \stopMPpage%
    }
    % pause at the first line, after the second one is made to coincide it
    \dorecurse{3*30}{%
        \startMPpage
            mytext ;
            lineone ;
            jasper_render_segments ;
        \stopMPpage
    }
    \startbuffer[BufA]
        \startframed[myframed]
            {\bf AXIOM 1} Given any two points in space, one and 
            only one line can be drawn through these points.\par
            It follows from AXIOM 1 that if any two points of a 
            line coincide with two points of some other line, 
            then these lines coincide in all their points.
            (Otherwise it would be possible to draw two 
            distinct lines through a given pair of points).\par 
            \color[white]{Therefore, two distinct lines may intersect only at 
            one point.\par 
            Let us emphasize that AXIOM 1 asserts not only the 
            existence of a line through a pair of points but also 
            the fact that such a line can be drawn.\par 
            Thus AXIOM 1 implies the existence of an instrument 
            that allows us to draw a straight line through any 
            pair of points. Such an instrument is called a 
            straightedge.}
        \stopframed
    \stopbuffer
    % (Otherwise it would be possible to draw two 
    % distinct lines through a given pair of points).
    \dorecurse{10*30}{%
        \startMPpage
            mytext ;
            lineone ;
            jasper_render_segments ;
        \stopMPpage
    }
    \startbuffer[BufA]
        \startframed[myframed]
            {\bf AXIOM 1} Given any two points in space, one and 
            only one line can be drawn through these points.\par
            It follows from AXIOM 1 that if any two points of a 
            line coincide with two points of some other line, 
            then these lines coincide in all their points.
            (Otherwise it would be possible to draw two 
            distinct lines through a given pair of points).\par 
            Therefore, two distinct lines may intersect only at 
            one point.\par 
            \color[white]{Let us emphasize that AXIOM 1 asserts not only the 
            existence of a line through a pair of points but also 
            the fact that such a line can be drawn.\par 
            Thus AXIOM 1 implies the existence of an instrument 
            that allows us to draw a straight line through any 
            pair of points. Such an instrument is called a 
            straightedge.}
        \stopframed
    \stopbuffer
    % Therefore, two distinct lines may intersect only at 
    % one point.
    \dorecurse{10*30}{%
        \startMPpage
            mytext ;
            lineone ;
            jasper_render_segments ;
        \stopMPpage
    }
    \def\myT{
        matrix_multiply(
            inverse(translate(cx,cy-5,0))
            ,matrix_multiply(
                zrotation(pi/6)
                ,translate(cx,cy-5,0)
            )
        )
    }%
    % show an example rotation
    \dorecurse{3*30}{%
        \startMPpage
            mytext ;
            lineone ;
            
            jasper_point[
                fx="cx"
                ,fy="cy+5"
                ,transformation="
                    matrix_multiply(
                        inverse(translate(cx,cy-5,0))
                        ,matrix_multiply(
                            zrotation(pi/6*#1/90)
                            ,translate(cx,cy-5,0)
                        )
                    )
                "
            ] ;
            jasper_point[
                fx="cx"
                ,fy="cy-5"
                ,transformation="
                    matrix_multiply(
                        inverse(translate(cx,cy-5,0))
                        ,matrix_multiply(
                            zrotation(pi/6*#1/90)
                            ,translate(cx,cy-5,0)
                        )
                    )
                "
            ] ;
            jasper_curve[
                ustart="0"
                ,ustop="1"
                ,usamples="2"
                ,fx="cx"
                ,fy="cy+7-14*u"
                ,transformation="
                    matrix_multiply(
                        inverse(translate(cx,cy-5,0))
                        ,matrix_multiply(
                            zrotation(pi/6*#1/90)
                            ,translate(cx,cy-5,0)
                        )
                    )
                "
            ] ;
            jasper_render_segments ;
        \stopMPpage%
    }
    % undo the rotation
    \dorecurse{3*30}{%
        \startMPpage
            mytext ;
            lineone ;
            
            jasper_point[
                fx="cx"
                ,fy="cy+5"
                ,transformation="
                    matrix_multiply(
                        inverse(translate(cx,cy-5,0))
                        ,matrix_multiply(
                            zrotation(pi/6*(1-#1/90))
                            ,translate(cx,cy-5,0)
                        )
                    )
                "
            ] ;
            jasper_point[
                fx="cx"
                ,fy="cy-5"
                ,transformation="
                    matrix_multiply(
                        inverse(translate(cx,cy-5,0))
                        ,matrix_multiply(
                            zrotation(pi/6*(1-#1/90))
                            ,translate(cx,cy-5,0)
                        )
                    )
                "
            ] ;
            jasper_curve[
                ustart="0"
                ,ustop="1"
                ,usamples="2"
                ,fx="cx"
                ,fy="cy+7-14*u"
                ,transformation="
                    matrix_multiply(
                        inverse(translate(cx,cy-5,0))
                        ,matrix_multiply(
                            zrotation(pi/6*(1-#1/90))
                            ,translate(cx,cy-5,0)
                        )
                    )
                "
            ] ;
            jasper_render_segments ;
        \stopMPpage%
    }
    % pause after completing both rotations
    \dorecurse{4*30}{%
        \startMPpage
            mytext ;
            lineone ;
            jasper_render_segments ;
        \stopMPpage
    }
    \startbuffer[BufA]
        \startframed[myframed]
            {\bf AXIOM 1} Given any two points in space, one and 
            only one line can be drawn through these points.\par
            It follows from AXIOM 1 that if any two points of a 
            line coincide with two points of some other line, 
            then these lines coincide in all their points.
            (Otherwise it would be possible to draw two 
            distinct lines through a given pair of points).\par 
            Therefore, two distinct lines may intersect only at 
            one point.\par 
            Let us emphasize that AXIOM 1 asserts not only the 
            existence of a line through a pair of points but also 
            the fact that such a line can be drawn.\par 
            \color[white]{Thus AXIOM 1 implies the existence of an instrument 
            that allows us to draw a straight line through any 
            pair of points. Such an instrument is called a 
            straightedge.}
        \stopframed
    \stopbuffer
    % Let us emphasize that AXIOM 1 asserts not only the 
    % existence of a line through a pair of points but also 
    % the fact that such a line can be drawn.
    \dorecurse{10*30}{%
        \startMPpage
            mytext ;
            lineone ;
            jasper_render_segments ;
        \stopMPpage
    }
    \startbuffer[BufA]
        \startframed[myframed]
            {\bf AXIOM 1} Given any two points in space, one and 
            only one line can be drawn through these points.\par
            It follows from AXIOM 1 that if any two points of a 
            line coincide with two points of some other line, 
            then these lines coincide in all their points.
            (Otherwise it would be possible to draw two 
            distinct lines through a given pair of points).\par 
            Therefore, two distinct lines may intersect only at 
            one point.\par 
            Let us emphasize that AXIOM 1 asserts not only the 
            existence of a line through a pair of points but also 
            the fact that such a line can be drawn.\par 
            Thus AXIOM 1 implies the existence of an instrument 
            that allows us to draw a straight line through any 
            pair of points. \color[white]{Such an instrument is called a 
            straightedge.}
        \stopframed
    \stopbuffer
    % Thus AXIOM 1 implies the existence of an instrument 
    % that allows us to draw a straight line through any 
    % pair of points.
    \dorecurse{10*30}{%
        \startMPpage
            mytext ;
            lineone ;
            jasper_render_segments ;
        \stopMPpage
    }
    \startbuffer[BufA]
        \startframed[myframed]
            {\bf AXIOM 1} Given any two points in space, one and 
            only one line can be drawn through these points.\par
            It follows from AXIOM 1 that if any two points of a 
            line coincide with two points of some other line, 
            then these lines coincide in all their points.
            (Otherwise it would be possible to draw two 
            distinct lines through a given pair of points).\par 
            Therefore, two distinct lines may intersect only at 
            one point.\par 
            Let us emphasize that AXIOM 1 asserts not only the 
            existence of a line through a pair of points but also 
            the fact that such a line can be drawn.\par 
            Thus AXIOM 1 implies the existence of an instrument 
            that allows us to draw a straight line through any 
            pair of points. Such an instrument is called a 
            straightedge.
        \stopframed
    \stopbuffer
    % Such an instrument is called a 
    % straightedge.
    \dorecurse{15*30}{%
        \startMPpage
            mytext ;
            lineone ;
            jasper_render_segments ;
        \stopMPpage
    }
\stoptext
